% Converted from README.md
\subsubsection{Memcheck}\label{memcheck}

Memcheck detektuje memorijske greške korisničkog programa. Kako ne vrši analizu izvornog koda već mašinskog, Memcheck ima mogućnost analize programa pisanom u bilo kom programskom jeziku. Za programe pisane u jezicima C i C++ detektuje sledeće probleme:

\begin{itemize}
\tightlist
\item
  Korišćenje nedefinisanih vrednosti, vrednosti koje nisu inicijalizovane ili koje su izvedene od drugih nedefinisanih. vrednosti. Problem se detektuje tek kada su upotrebljene.
\item
  Čitanje ili pisanje u nedopuštenu memoriju na hipu, steku, bilo da je potkoračenje ili prekoračenje dozvoljene memorije ili pristupanje već oslobodenoj memoriji.
\item
  Neispravno oslobadanje memorije na hipu, npr. duplo oslobadanje memorije na hipu ili neupareno korišćenje funkcija \texttt{malloc/new/new{[}{]}} i \texttt{free/delete/delete{[}{]}}.
\item
  Poklapanje argumenata \texttt{src} i \texttt{dest} funkcije \texttt{memcpy} i njoj sličnim.
\item
  Prosledivanje loših vrednosti za veličinu memorijskog prostora funkcijama za alokaciju memorije, npr. negativnih.
\item
  Curenje memorije, npr. gubitak pokazivača na alociran prostor.
\end{itemize}

\subimport{01_uninitialized/}{index.tex}

\subimport{02_undefined/}{index.tex}

\subimport{03_malloc/}{index.tex}

\subimport{04_new_delete/}{index.tex}

\subimport{05_overlap/}{index.tex}

\subimport{06_fishy_arguments/}{index.tex}

\subimport{07_misc/}{index.tex}
